\begin{afterword}
\summaryhead

The post continues ``Lotteries: A Waste of Hope,'' published the same day. Defenders had called the lottery a cheap way to buy a daydream, and a commenter had written that a ticket turns zero chance of wealth into epsilon. The author replies that the difference between zero and epsilon is only of order epsilon.

Then comes a satire. If what the lottery sells is epsilon hope, a better product would be the New Improved Lottery: one dollar buys a chance at a prize paid at a random moment, about once every five years. The buyer could stare at the phone at work, fill a shopping cart and put things back if no call comes, and watch live odds on a phone screen, until this daydream crowded out every other. Governments do not sell this, the post says, because they want people to pay every week. So if the lottery is a service, it is an overpriced one, charged to the poorest, and citizens should demand the new version.

\responsehead

The satire makes one point well. If what a ticket buys is hope, current lotteries charge a great deal for it. At the 2007 price of a dollar a play and two draws a week, a player who bought a ticket for every draw of one game would pay about \$520 over the five years that one dollar buys in the New Improved Lottery. The post draws that conclusion, and on its premise the arithmetic supports it.

The other arguments are weaker. The reply to the commenter, that zero and epsilon differ by only epsilon, is true of the probabilities. The commenter's claim was about the dream: that any real chance lets it ``bridge that gap.'' Whether a daydream's worth should shrink with its probability is the question at issue, and the reply assumes that it should. The same comment also said that ``nobody can measure the value of entertainment or distant hope for another.'' The post quotes only the first part of the comment, and its satire answers the price of hope, not whether its value can be judged from outside.

The satire's picture of harm, a buyer glued to a screen instead of working or studying, is an exaggerated version of the previous post's claim that lottery dreams displace better ones. That claim was offered there with ``maybe'' and ``might,'' and the satire adds no evidence for it.

One factual point bears on the ending. The post says governments ``don't offer this simple and obvious service.'' Since 1956 the United Kingdom has sold Premium Bonds, which enter one purchase in a monthly prize draw for as long as it is held. They are not the post's design, since the draws are monthly and the stake can be reclaimed, but a government does sell long-lasting chances for a single purchase.

In short: a sound point about the price of hope, a reply to the commenter that assumes the point in dispute, and a satire whose harms restate a claim the previous post made without evidence.
\end{afterword}
